\documentclass[conference]{IEEEtran}

\usepackage{booktabs}
\usepackage{cite}
\usepackage{graphicx}
\usepackage[hidelinks]{hyperref}
\usepackage{listings}
\usepackage{url}
\usepackage{xspace}

\newcommand{\tool}{SafeMAP\xspace}
\newif\ifanonymous
\anonymousfalse % Change only if the eventual venue requires anonymous review.

\begin{document}

\title{SafeMAP: Analysis-Guided Planning and Validation for Safe C-to-Rust Translation}

\ifanonymous
\author{\IEEEauthorblockN{Anonymous Submission}}
\else
\author{
\IEEEauthorblockN{1\textsuperscript{st} Jacob Sebastian Cyril}
\IEEEauthorblockA{\textit{Department of Computer Science}\\
\textit{CHRIST (Deemed to be University)}\\
Bangalore, India\\
jacob.sebastian@mca.christuniversity.in}
\and
\IEEEauthorblockN{2\textsuperscript{nd} Benita Jaison}
\IEEEauthorblockA{\textit{Department of Computer Science}\\
\textit{CHRIST (Deemed to be University)}\\
Bangalore, India\\
email address or ORCID}
\and
\IEEEauthorblockN{3\textsuperscript{rd} New Begin M.}
\IEEEauthorblockA{\textit{Department of Computer Science}\\
\textit{CHRIST (Deemed to be University)}\\
Bangalore, India\\
email address or ORCID}
}
\fi

\maketitle

\begin{abstract}
Automatic C-to-Rust translators often preserve C pointer conventions and emit
unsafe Rust. \tool takes a different approach. It records each migration
decision, from candidate classification through behavioral testing, in an
auditable plan. Its analyzer collects pointer and dependency facts. A
deterministic synthesizer emits Rust only when a safe rule matches exactly.
Accepted code must compile without unsafe operations, expose no raw pointers in
its public API, and match the original C on recorded randomized inputs. The C
oracle runs under AddressSanitizer and UndefinedBehaviorSanitizer. Six helpers
from a 32-function LLVM development corpus pass every check; the previous rule
set supported one. Each helper also passes 1,000 differential cases. We specify
a frozen held-out study on three independently maintained C libraries. The
protocol uses outcome-blind human labels, records analyzer provenance, compares
against direct LLM translation to safe Rust, and tests individual components
through ablation. At manuscript freeze, held-out outcomes and reviewer
agreement replace this protocol-only statement. Development results test the
mechanism but do not measure generalization.
\end{abstract}

\begin{IEEEkeywords}
C-to-Rust migration, safe Rust, static analysis, program translation, software
maintenance, memory safety
\end{IEEEkeywords}

\section{Introduction}

C remains common in systems software. Its memory model leaves programmers
responsible for preventing buffer overflows, use-after-free errors, null
dereferences, and data races. Safe Rust checks ownership and borrowing at
compile time while retaining low-level control. Translation is difficult
because many C idioms have no direct safe-Rust equivalent. Raw pointers, output
parameters, manual allocation, nullable values, and integer error codes require
changes in representation and API design.

The established C2Rust translator preserves C behavior by generating Rust that
often uses raw pointers and unsafe operations \cite{c2rust}. Its output is a
useful reference, although it does not provide the guarantees of safe Rust.
Prior work uses static analysis, refactoring, ownership-guided rewriting,
transpiler redesign, and LLM-based rewriting to remove unsafe code after
translation
\cite{ling2022trust,emre2021safer,hong2023static,zhang2023ownership,
nitin2026c2saferrust,mohapatra2026cp2srust,hong2025cacm}. Other studies ask
whether LLMs can translate directly to idiomatic Rust and repair failures with
compiler or test feedback
\cite{hong2025typemigrating,zhou2025c2rusttv,eniser2025fluorine,
valenzuela2026llmcompile,zhang2026rustflow}. A third line
targets safe Rust directly for structured subsets of C \cite{fromherz2026scylla}.

\tool also targets safe Rust, but it is an experimental migration pipeline, not
a general C compiler. Its acceptance rule is strict. Reducing unsafe code is
insufficient, as is compilation with unsafe blocks or a C-style raw-pointer
API. \tool accepts a unit only when the generated Rust satisfies the project's
safe-code policy and passes every applicable check. If no supported safe path
exists, the tool rejects the unit and records why.

We ask:

\begin{itemize}
  \item \textbf{RQ1:} For a restricted, statically identifiable subset of C
  idioms, how often can \tool synthesize fully safe Rust?
  \item \textbf{RQ2:} Which C idioms are currently handled reliably by
  deterministic safe synthesis, and which idioms remain limiting?
  \item \textbf{RQ3:} How does \tool compare with direct LLM safe translation
  and raw C2Rust under one acceptance policy?
  \item \textbf{RQ4:} Which individual planning and validation components
  contribute to accepted translation?
  \item \textbf{RQ5:} Which artifacts are required to audit and reproduce a
  safe-first migration decision?
\end{itemize}

The evaluation keeps development and test data separate. Development data
contains 40 authored microbenchmarks, six authored regression modules, and ten
LLVM programs whose results were already known. The frozen test corpus retains
every function definition in preselected production source files from pinned
revisions of inih, cJSON, and libcsv. The files stay sealed until the
implementation is frozen.

This paper contributes:

\begin{itemize}
  \item a C-to-Rust pipeline that accepts code only when it contains no unsafe
  Rust and exposes no raw pointers through its public API;
  \item a migration plan that records detected C idioms, proposed Rust APIs,
  eligibility decisions, and required checks;
  \item a deterministic synthesizer for supported C idioms, connected to
  compilation, source-policy checks, differential testing, and recorded
  sanitizer status; and
  \item a pre-specified function-level study with independent eligibility
  labels, analyzer provenance, replayable test cases, component ablations, and
  a direct safe-output baseline.
\end{itemize}

\section{Background and related work}

\subsection{C2Rust and safer-Rust refactoring}

C2Rust translates C through a mostly syntactic process and is often used as a
baseline or a migration starting point \cite{c2rust}. It preserves many C
constructs with raw pointers and unsafe operations, so much of its output falls
outside Rust's safe subset. Emre et al. found that raw pointers cause much of
this unsafety \cite{emre2021safer}. Hong used static analysis to replace
selected unsafe constructs in C2Rust output, including lock APIs and output
parameters \cite{hong2023static}.

C2SaferRust combines C2Rust, static analysis, and LLM rewriting to reduce unsafe
code in translated real-world programs \cite{nitin2026c2saferrust}. Crown
infers ownership models for C pointers and converts many of them to safe Rust
types \cite{zhang2023ownership}. CRustS and Cp2SRust use different
source-to-source transformations and transpiler designs
\cite{ling2022trust,mohapatra2026cp2srust}. Industry studies and surveys report
that automatic translation still needs substantial manual work and testing
\cite{hong2025cacm,viirola2026feasible}. \tool uses C2Rust as a reference lane.
Its primary lane generates final safe Rust directly for eligible units.
Reductions in unsafe code remain diagnostic measurements; they count as
migrations only when the final output passes the stated acceptance policy.

\subsection{LLM-based C-to-Rust translation}

LLMs can produce idiomatic Rust, although correctness and type migration remain
open problems. Hong and Ryu studied an LLM translator supplied with candidate
signatures, caller and callee context, and compiler feedback
\cite{hong2025typemigrating}. C2RustTV uses LLMs for production and test-code
translation, then checks the result with generated tests and repair loops
\cite{zhou2025c2rusttv}. Flourine evaluates LLM translation through
cross-language differential fuzzing and finds lower success on more complex
real-world samples \cite{eniser2025fluorine}. RustFlow studies prompting and
iterative repair \cite{zhang2026rustflow}. Valenzuela et al. examine compiler
errors in direct C-to-safe-Rust translation and in LLM repair of C2Rust output
\cite{valenzuela2026llmcompile}. RustMap combines program analysis and LLM
support at project scale \cite{cai2025rustmap}.

\subsection{Safe-Rust translation for C subsets}

Scylla translates an applicative C subset to safe Rust and shows how source
structure can make safe generation possible \cite{fromherz2026scylla}. \tool
likewise makes no claim that arbitrary C can be migrated automatically. It is a
small research prototype built around migration plans, recorded rejections, and
reproducible acceptance measurements.

\subsection{Validation and undefined behavior}

Compilation alone cannot establish a safe migration. \tool records compilation
status, remaining unsafe constructs, raw pointers in public APIs, differential
test results, C sanitizer results, and Miri status. Miri can detect undefined
behavior during Rust execution \cite{jung2026miri}. ASan and UBSan check the
executed C cases before \tool uses them as behavioral oracles
\cite{serebryany2012asan,clangubsan}. Miri is disabled in the frozen experiment
because the generated crates and installed toolchain do not provide a uniform
applicable oracle; the artifact reports it as skipped rather than passed.

\section{SafeMAP design}

\subsection{Outcome labels}

\tool records candidate classification and implementation support separately.
The candidate label is \texttt{candidate\_safe},
\texttt{manual\_refactor\_required}, \texttt{unsafe\_required}, or
\texttt{unknown}. A candidate-safe label means that screening found no known
obstacle. It does not prove that a safe translation exists. The separate
\texttt{implemented\_support} field is true only when an exact synthesis rule
matches. Later fields record generation, compilation, policy compliance, and
behavioral test results. Every decision also identifies whether its facts came
from libclang or the regular-expression fallback.

\subsection{Acceptance policy}
\label{sec:acceptance-policy}

\tool reports policy and behavioral acceptance separately. A policy-safe unit
must meet these conditions:

\begin{enumerate}
  \item The final Rust code compiles with \texttt{\#![forbid(unsafe\_code)]}.
  \item The final Rust contains no \texttt{unsafe} blocks, \texttt{unsafe fn},
  \texttt{unsafe impl}, or \texttt{extern "C"} blocks.
  \item The public API does not expose raw pointers such as \texttt{*const T}
  or \texttt{*mut T}.
  \item The planned function or module unit is present in the final Rust output.
\end{enumerate}

A behaviorally validated unit is policy-safe and passes an applicable
C-versus-Rust oracle. A skipped or inapplicable check does not count as a pass.
Cargo test and Clippy must pass when enabled. The frozen run records Miri as
disabled and never treats that status as a pass. ASan and UBSan monitor the C oracle; a
sanitizer finding disqualifies that execution as a correctness oracle.
Partially migrated C2Rust output may still aid manual work, but \tool does not
accept it unless it passes the same policy.

\subsection{Pipeline}

\tool processes each C input in seven stages.
Figure~\ref{fig:safemap-pipeline} shows both execution lanes and the policy they
share:

\begin{enumerate}
  \item Locate source files and recover compile information when possible.
  \item Identify functions, signatures, pointer use, allocation patterns,
  unsupported constructs, and dependency structure.
  \item Assign the screening taxonomy without consulting implemented synthesis
  rules.
  \item Classify migration idioms such as output
  parameters, pointer-length arrays, nullable pointers, C strings, and
  return-code-plus-output patterns.
  \item Build a plan that records the target Rust signature, ownership model,
  validation requirements, and reasons for any rejection.
  \item Generate Rust when a deterministic rule supports the planned pattern.
  \item Compile the result, check remaining unsafe code, count raw pointers, run
  differential tests, check ASan/UBSan on C, record Miri status, and export
  function-level measurements.
\end{enumerate}

\subsection{Worked plan and rule selection}

Consider the LLVM development helper:
\begin{lstlisting}[language=C,basicstyle=\ttfamily\scriptsize]
double loop(float *x, float *y, long n) {
  double acc = 0.0;
  for (long i=0; i<n; ++i)
    acc += (double)x[i] * (double)y[i];
  return acc;
}
\end{lstlisting}
libclang finds two indexed pointers paired with \texttt{n}. It finds no indexed
writes, calls, or dependencies. The classifier returns
\texttt{candidate\_safe}. The plan maps each pointer-length pair to a shared
slice, and the \texttt{slice\_dot\_product} rule provides
\texttt{implemented\_support}. Since \texttt{loop} is a Rust keyword, the
generated function uses a raw identifier:
\begin{lstlisting}[basicstyle=\ttfamily\scriptsize]
pub fn r#loop(x: &[f32], y: &[f32]) -> f64 {
  x.iter().zip(y.iter())
   .map(|(&a,&b)| (a as f64)*(b as f64))
   .fold(0.0_f64, |s,v| s + v)
}
\end{lstlisting}
The initial value \texttt{+0.0} is necessary. Differential testing showed that
a generic floating-point \texttt{sum()} could disagree with C on the sign bit
of zero. The generated helper matches C bit for bit on 1,000 recorded cases and
passes the Cargo, Clippy, ASan, and UBSan checks. LLVM's \texttt{doTest}, by
comparison, calls allocation and memory APIs outside the supported subset. The
classifier labels it \texttt{unsafe\_required}, so synthesis does not run.

\begin{lstlisting}[basicstyle=\ttfamily\scriptsize]
facts <- analyze(function)
candidate <- classify(facts)
unit <- SCC(call_graph, function)
plan <- safe_signature(facts) + idiom_facts
if candidate != candidate_safe:
    support <- not_applicable
else if exact_rule_matches(function, plan):
    support <- implemented_support
else: support <- not_implemented
generate only with implemented_support
report policy and behavior separately
\end{lstlisting}

\begin{figure*}[t]
  \centering
  \includegraphics[width=\textwidth]{figures/safemap_conversion_pipeline.pdf}
  \caption{\tool conversion and acceptance pipeline. Static analysis records
  eligibility separately from synthesis support. The primary lane generates
  Rust from a migration plan; C2Rust runs in a separate reference lane. The
  same policy applies to both outputs. Reports retain unsupported cases,
  missing synthesis rules, and failed checks as distinct outcomes.}
  \label{fig:safemap-pipeline}
\end{figure*}

\subsection{Supported idioms}

The deterministic synthesizer currently handles:
integer and floating scalar expressions, integer boolean and bit-reversal
idioms, pointer-length and fixed-size arrays, scalar structs and field access,
dependency-ordered internal helper calls, simple entry-point composition,
mutable scalar pointer updates, output
parameters, multiple output parameters, return-code plus output-parameter to
\texttt{Result}, nullable pointers, reductions including two-slice floating
dot products, simple C string lengths, simple allocation-to-\texttt{Vec}
patterns, and single owned allocations to \texttt{Box}. It does not support
unions, inline assembly, volatile memory access, unresolved aliasing, complex
macros, custom allocators, function pointers, or pointer-integer casts.

\subsection{Baselines}

\tool compares its primary lane with two baselines. Raw C2Rust measures how much
unmodified output passes the policy. The second baseline asks a fixed LLM to
translate C directly to safe Rust without \tool facts or plans. The frozen
baseline is Google's \texttt{gemini-3.5-flash-lite}, accessed through its
OpenAI-compatible API at temperature zero. The experiment keeps each prompt
and response along with model settings, token counts, retries, and failures. It
makes one request per project and performs no automatic retry; quota or provider
failures remain in the denominator. Five ablations remove one component at a time: pointer-role
classification, safe-signature generation, dependency grouping, idiom plans,
or validation feedback. Every lane uses the same source units and behavioral
checks.

\section{Implementation}

\tool is a Python package. Separate modules handle ingestion, static analysis,
migration planning, translation, repair, validation, and measurement. For a
supported idiom, the deterministic synthesizer emits Rust under
\texttt{\#![forbid(unsafe\_code)]}. The validation layer derives scalar and
array harnesses from function signatures. It saves seeds and concrete C/Rust
inputs and outputs, then records the result of each compiler, policy,
differential, sanitizer, Clippy, and recorded Miri status. The benchmark runner writes
CSV, Markdown, and LaTeX tables. One evaluation command combines \tool,
case-study, external-corpus, C2Rust, and ablation results. Its CSV files also
record function roles, complete-project outcomes, outcome stages, and the
status of each check.

The analyzer uses libclang when available. Otherwise it uses a pattern-based
parser and records why the fallback was needed. Dependency analysis groups
mutually dependent functions into one unit. Classification and planning finish
before synthesis begins. The deterministic lane needs neither C2Rust nor a
network model service.

\section{Evaluation methodology}

\subsection{Datasets}

The development corpus contains 40 authored microbenchmarks, six authored
modules, and ten LLVM test-suite programs used during earlier experiments. The
LLVM sources are pinned at commit \texttt{6cdc54e} \cite{llvmtestsuite}. They
contain 589 lines of C and 32 functions. libclang analyzed every function.
Their outcomes informed rule development, so they provide no held-out evidence.

The sealed test corpus uses pinned production sources from inih r62, cJSON
v1.7.19, and libcsv commit \texttt{b1d5212}. A documented selection rule keeps
every function in named production \texttt{.c} files without consulting \tool
output. Before execution, the experiment records repository URLs, revisions,
source hashes, licenses, and physical line counts. A materialization script
blocks access until the implementation has a frozen Git commit. This keeps test
sources out of rule development.

Two reviewers label each held-out function without seeing the \tool decision.
They then resolve disagreements. The labels use the candidate taxonomy.
\texttt{unknown} cases are reported separately and excluded from binary
precision and recall. The study reports raw agreement, Cohen's $\kappa$,
disagreements by construct, and separate results for libclang and the fallback
backend.

\subsection{Metrics}

The primary effectiveness measure is the number of behaviorally validated,
policy-safe units divided by all retained held-out units. The report also gives
the count at each earlier stage: candidate-safe, supported, generated,
compiling, and policy-safe. Results are separated by declared targets, helpers,
entry points, other functions, and complete modules. Analyzer measurements
include eligibility precision and error counts. Corpus tables report project
and function counts, lines of code, cyclomatic complexity, pointer density, and
construct frequencies.

\subsection{Execution snapshot}

Before unsealing, the implementation and configuration are committed. A freeze
record binds the experiment to that Git revision. The held-out test then runs
once. Artifacts record source and configuration hashes, tool versions, analyzer
backend, commands, generated code, differential inputs and outputs, random
seeds, sanitizer status, timing, and failures. The direct-LLM baseline fixes
the model and sampling settings before this run. The implementation-freeze
record stores the exact pre-freeze test command and result; the test count is
not a held-out outcome.

\section{Results}

\subsection{Development evidence}

The LLVM development corpus contains 32 functions from ten projects. After two
analysis fixes and rule changes based only on development data, 23 functions
receive \texttt{candidate\_safe} and six receive
\texttt{implemented\_support}. The six generated functions compile, satisfy
the policy, and pass behavioral testing. They are \texttt{sqr}, \texttt{fade},
\texttt{lerp}, the 32-bit and 64-bit reversal helpers, and the two-slice
floating-point reduction. Before these changes, only \texttt{sqr} passed. The
6/32 result measures development progress and exercises the reporting
pipeline. It cannot support a generalization claim because the same corpus
guided the rule changes.

\begin{table}[t]
\centering
\caption{Development-only outcome cascade. ``Cand.'' denotes
\texttt{candidate\_safe}; ``Behav.'' denotes behaviorally validated
policy-safe output.}
\label{tab:development-cascade}
\resizebox{\columnwidth}{!}{%
\begin{tabular}{lrrrrrrr}
\toprule
Corpus & Total & Cand. & Support & Gen. & Compile & Policy & Behav. \\
\midrule
LLVM programs & 32 & 23 & 6 & 6 & 6 & 6 & 6 \\
Microbenchmarks & 81 & 76 & 36 & 36 & 36 & 36 & 36 \\
Authored modules & 25 & 25 & 19 & 19 & 19 & 19 & 17 \\
\bottomrule
\end{tabular}}
\end{table}

Two bugs show why the report keeps intermediate outcomes. The pointer screen
first treated indexed reads as writes, which produced a false mutable-alias
conflict for the floating-point reduction. It also parsed \texttt{>>} as a
comparison and gave bit reversal a Boolean return type. Both bugs now have
regression tests. Exact differential comparison found the separate signed-zero
error described earlier.

On the 40 authored development programs, the full configuration validates 36
of 81 units, equal to 36 of 76 candidate-safe units, and 35 of 40 declared
targets. Without pointer roles, the counts fall to 16/81 units and 15/40
targets. Without safe signatures or idiom plans, they fall to 6/81 and 5/40.
Removing dependency grouping or validation feedback has no effect on this
suite. These development ablations are diagnostic. The zero changes for the
last two components show that the authored programs do not test their proposed
benefits, so no causal conclusion follows.

A sixth authored module tests fixed arrays, a scalar struct, field access, two
internal calls, five functions in one crate, and a small \texttt{main}. All
five functions compile, pass policy checks, and pass the generated differential
and sanitizer harness. Across six authored modules, the full system validates
17/25 units. Removing dependency grouping reduces the result to 15/25 because
the two composed callers lose support. The module was written to test this
mechanism, so the result is development evidence.

\begin{table}[t]
\centering
\caption{Behaviorally validated units in component ablations. Denominators
include all functions or dependency units in each development corpus.}
\label{tab:development-ablations}
\begin{tabular}{lrr}
\toprule
Configuration & Microbenchmarks & Modules \\
\midrule
Full system & 36/81 & 17/25 \\
Without pointer roles & 16/81 & 3/25 \\
Without safe signatures & 6/81 & 0/25 \\
Without dependency grouping & 36/81 & 15/25 \\
Without idiom plans & 6/81 & 1/25 \\
Without validation feedback & 36/81 & 17/25 \\
\bottomrule
\end{tabular}
\end{table}

\subsection{Held-out evaluation status}

The held-out corpus has not been materialized or run for this manuscript.
There is therefore no effectiveness estimate for RQ1--RQ4. We also do not
combine authored and LLVM development results into a substitute test rate.
After the implementation freeze, the artifact will generate final tables for
corpus composition, label agreement, eligibility errors, outcomes by function
role, complete modules, both baselines, and the five component ablations. Those
results must replace this section before submission.

\section{Discussion}

\subsection{How the acceptance rule changes comparison}

Code that compiles while exposing raw pointers through its public API has not
reached safe Rust. The same policy therefore applies to every experimental
lane. This prevents raw C2Rust output from being counted beside policy-safe
output merely because both compile. C2Rust remains useful as a reference, and
its output passes only when it meets the shared policy.

\subsection{Isolating planning components}

The synthesizer requires a plan, so removing every form of guidance guarantees
failure. The main ablation instead removes one input per run: pointer roles,
safe signatures, dependency grouping, idiom plans, or validation feedback. A
no-guidance run remains as a dependency check. Source units and validators stay
fixed across runs.

\subsection{Recorded rejection}

Unsupported constructs such as unions, volatile access, inline assembly, and
function pointers occur in real C programs. Partial output can hide their
effect on a migration. \tool records each unsupported case as an outcome. The
record shows where automatic processing stopped and where manual review must
begin.

\subsection{Eligibility and synthesis coverage}

The authored suites contain deliberately supported patterns, so they measure
regression coverage. In the LLVM development set, 17 of 23 candidate-safe
functions still lack a synthesis rule. Candidate labels and implementation
support must therefore remain separate. Held-out human labels are needed to
tell how much of the gap comes from classification error and how much comes
from missing rules.

\section{Threats to validity}

\textbf{Benchmark representativeness.} The authored and LLVM corpora are
development data. The held-out set improves project diversity but contains
three small C libraries instead of large applications. We report size,
complexity, pointer use, and construct frequencies so readers can judge the
scope. The results do not cover arbitrary C.

\textbf{Implementation maturity.} \tool is a research prototype. The analyzer
and synthesizer use patterns and do not implement full C semantics. The tool
reports unsupported constructs, and its supported subset remains small.

\textbf{Differential testing limits.} Differential testing over bounded random
inputs is evidence, not equivalence proof. Generated harnesses cover supported
scalar and array shapes. They do not cover every heap, callback, or stateful
API. Saved seeds and outputs allow replay. ASan and UBSan reduce the chance of
using an undefined C execution as an oracle, but they cannot say whether
unexecuted inputs are defined.

\textbf{Label subjectivity.} Candidate safety sometimes requires judgment
about whether an API redesign is automatic or manual. Two outcome-blind labels,
retained disagreements, adjudication, confidence scores, and $\kappa$ expose
that judgment in the report.

\textbf{Toolchain sensitivity.} C2Rust, Clang, Rust, Clippy, and Miri can
behave differently across host versions and configurations. The artifact
records skipped and failed stages. Reproduction across machines will need
tighter toolchain pinning.

\textbf{Metric interpretation.} Units are functions or dependency groups, not
lines of code or whole-project migration percentages. Analyzer eligibility may
exceed implemented synthesis support. Dataset size and complexity describe the
corpus but do not remove this measurement limit.

\section{Future work}

Current rules handle fixed arrays, scalar structs and fields, internal calls,
small entry-point compositions, and multi-function crates. Planned work covers
stateful APIs, callbacks, nested aggregates, unions with safe tagged
representations, and aliases that require flow-sensitive borrowing analysis.
Test generation also needs structured inputs, buffers with semantic length
constraints, and preconditions derived from metadata. If held-out results lead
to rule changes, the next generalization study will need a new frozen corpus.

\section{Conclusion}

\tool records each stage of C-to-Rust migration instead of reducing the result
to compilation success. Classification, synthesis support, policy checks, and
behavioral tests have separate outcomes. A missing test harness cannot count as
a behavioral pass. Six functions in the LLVM development corpus now pass every
stage, among them a floating-point reduction and two bit-reversal helpers. The
held-out protocol is prepared for three independent projects and includes
outcome-blind labels, replayable differential cases, C sanitizers, a direct
safe-output baseline, and component ablations. Until that study runs, the
results establish only that the implemented rules and measurement pipeline work
on development data.

\bibliographystyle{IEEEtran}
\bibliography{references}

\end{document}
